\documentclass[a4paper,12pt]{article}

\usepackage[a4paper,margin=1in]{geometry}
\usepackage{enumitem}
\usepackage{listings}
\usepackage{xcolor}
\usepackage{hyperref}
\usepackage{graphicx}
\usepackage{float}
\usepackage{array}
\usepackage{longtable}

\setlength{\parindent}{0pt}
\setlength{\parskip}{0.6em}
\setlist[itemize]{itemsep=0.35em}

\lstdefinestyle{CppStyle}{
    language=C++,
    basicstyle=\ttfamily\footnotesize,
    keywordstyle=\color{blue},
    commentstyle=\color{green!40!black},
    numbers=left,
    numberstyle=\tiny,
    breaklines=true,
    frame=single,
    columns=fullflexible,
    showstringspaces=false
}

\title{Project 3 Report: 2048 Game\\
\large CMPE 230 Systems Programming - Spring 2026}

\author{
Kemal Önal 
}

\date{\today}

\begin{document}

\maketitle

\section{Project Overview}

This project is an implementation of the classic 2048 sliding-tile puzzle game using the Qt Widgets framework. The goal of the game is to combine tiles with equal values until the player reaches the target tile value, which is 2048 by default. The game is played on an \(N \times M\) grid. Empty cells contain no tile, while occupied cells contain powers of two such as 2, 4, 8, 16, and so on.

The application is written in C++17 and built with CMake. The executable target is named \texttt{2048}, as required by the project specification. The implementation is divided into two main responsibilities. The \texttt{GameEngine} class stores the board and implements the game rules. The \texttt{MainWindow} class builds the Qt user interface, receives user input, displays scores and tiles, and keeps the visible board synchronized with the internal board.

The program starts in \texttt{main.cpp}. A \texttt{QApplication} object is created, a \texttt{MainWindow} object is shown, and Qt's event loop starts. From that point onward, the program reacts to events: keyboard presses, button clicks, timer timeouts, and message box choices.

The game supports three modes:

\begin{itemize}
    \item \textbf{Normal Mode:} The standard 2048 game. Reaching the target tile changes the game state to won.
    \item \textbf{Unlimited Mode:} The player can continue after reaching the target tile and create larger tiles.
    \item \textbf{Hard Mode:} If the player does not make a move within five seconds, the program automatically performs a random valid move. This mode has the same win and loss conditions as Normal Mode.
\end{itemize}

The user can restart the game at any time, and the game also supports unlimited Undo operations during active play. The best score is tracked across games within the same program session.

\section{User Interface Design}

The graphical interface is implemented in the \texttt{MainWindow} class. The UI is intentionally kept close to the visual style of 2048: a light background, a compact title area, score boxes, control buttons, and a colored tile board. The main layout is created programmatically with Qt layout classes rather than using a separate UI file.

The highest-level layout is a vertical page layout. Inside it, a centered content widget contains the game title, score row, mode and control buttons, and the board widget. The board widget contains a \texttt{QGridLayout}. Each cell in the grid is a \texttt{QLabel}; this makes it simple to update tile text, color, font, and size after every move.

\begin{figure}[H]
\centering
\includegraphics[width=0.85\textwidth]{230ss.png}
\caption{Screenshot of the 2048 user interface.}
\end{figure}

\subsection*{Layout Components}

The UI contains the following major components:

\begin{itemize}
    \item \textbf{Title area:} A large \texttt{2048} title and a short subtitle.
    \item \textbf{Score cards:} Two labels show the current score and the best score. They are styled as compact cards.
    \item \textbf{Mode buttons:} Normal, Unlimited, and Hard buttons switch game modes without resetting the board.
    \item \textbf{Control buttons:} Undo and Restart buttons provide explicit controls in addition to keyboard shortcuts.
    \item \textbf{Board container:} A styled outer widget contains the tile grid.
    \item \textbf{Tile labels:} Each tile is a \texttt{QLabel} whose text and background color are updated from the engine state.
\end{itemize}

The current UI setup creates a centered content widget and then places title, score, button, and board elements into it. The score labels are ordinary \texttt{QLabel} widgets, but their text is formatted with simple HTML to display a small title and a larger numeric value. This keeps the score display prominent without requiring a custom widget class.

\begin{lstlisting}[style=CppStyle,caption={Score label update}]
scoreLabel->setText(QString("<span style='font-size: 9px;'>SCORE</span><br>"
                            "<span style='font-size: 18px;'>%1</span>")
                        .arg(engine.getScore()));
bestScoreLabel->setText(QString("<span style='font-size: 9px;'>BEST</span><br>"
                                "<span style='font-size: 18px;'>%1</span>")
                            .arg(engine.getBestScore()));
\end{lstlisting}

\subsection*{Responsive Board Sizing}

The board size is configurable through constants. Because the board can be changed from the default 4 by 4 to another rectangular size, the tile labels should not have a fixed large size. The implementation therefore recalculates tile size according to the current window size and the configured row and column counts. This keeps tiles square and prevents large boards from becoming unusable.

The tile size is selected as the minimum of the size allowed by width and the size allowed by height. The value is also bounded so that the default 4 by 4 board does not become excessively large and large boards do not produce negative or unusably small tile sizes.

\begin{lstlisting}[style=CppStyle,caption={Dynamic square tile sizing}]
const int tileByWidth =
    (availableWidth - 2 * padding - spacing * (BOARD_M - 1)) / BOARD_M;
const int tileByHeight =
    (availableHeight - 2 * padding - spacing * (BOARD_N - 1)) / BOARD_N;
const int tileSize = std::max(26, std::min({tileByWidth, tileByHeight, 104}));

for (int r = 0; r < BOARD_N; ++r) {
    for (int c = 0; c < BOARD_M; ++c) {
        boardLabels[r][c]->setFixedSize(tileSize, tileSize);
        boardLabels[r][c]->setFont(tileFont);
    }
}
\end{lstlisting}

The \texttt{resizeEvent} function calls \texttt{updateTileSizes}, so the board is recalculated if the window or web viewport changes.

\section{Game State Representation}

The internal game state is represented in \texttt{GameEngine}. This class is independent from Qt widgets. It stores the board, score values, selected game mode, game state, and undo history. This separation keeps the game rules easier to reason about because movement and merging do not directly depend on UI code.

The most important configurable constants are located near the top of \texttt{gameengine.h}. These constants control the row count, column count, target tile value, and new tile probabilities.

\begin{lstlisting}[style=CppStyle,caption={Configurable game constants}]
constexpr int N = 4;      // Board row count
constexpr int M = 4;      // Board column count
constexpr int K = 2048;   // Target tile value
constexpr int P = 90;     // Chance of spawning a 2 (%)
constexpr int Q = 10;     // Chance of spawning a 4 (%)

static_assert(N > 0 && M > 0, "Board dimensions must be positive.");
static_assert(K > 0, "Target tile value must be positive.");
static_assert(P >= 0 && Q >= 0 && P + Q == 100,
              "P and Q must be non-negative and sum to 100.");
\end{lstlisting}

The code then maps these specification-style constants to names used by the rest of the implementation, such as \texttt{BOARD\_N}, \texttt{BOARD\_M}, \texttt{TARGET\_SCORE}, \texttt{CHANCE\_P}, and \texttt{CHANCE\_Q}. This keeps the rest of the code readable while still making the required constants easy to change near the top of the file.

The board itself is stored as a two-dimensional vector of integers:

\begin{lstlisting}[style=CppStyle,caption={Board and state fields in GameEngine}]
std::vector<std::vector<int>> board;
std::stack<GameStateData> history;

int score;
int bestScore;
GameMode currentMode;
GameState currentState;
\end{lstlisting}

An empty cell is represented by \texttt{0}. Non-empty cells contain tile values such as 2, 4, 8, 16, and larger powers of two. This representation is simple and efficient for scanning rows, columns, and neighbors.

The program uses three enums to keep the code readable: \texttt{Direction} for movement, \texttt{GameMode} for the selected mode, and \texttt{GameState} for whether the game is playing, won, or lost.

\subsection*{Restart State}

Restart creates a new board, resets the current score, returns the game state to \texttt{PLAYING}, clears the undo history, and spawns two initial tiles. The initial tiles are forced to value 2, as required by the specification.

\begin{lstlisting}[style=CppStyle,caption={Restart initializes a new game}]
void GameEngine::restart() {
    board.assign(BOARD_N, std::vector<int>(BOARD_M, 0));
    score = 0;
    currentState = GameState::PLAYING;

    while (!history.empty()) {
        history.pop();
    }

    spawnTile(2);
    spawnTile(2);
}
\end{lstlisting}

\subsection*{Undo State}

Undo is implemented with a stack. The stack stores only the board and the current score. It does not store the best score, because the best score is supposed to represent the highest score achieved during the session.

\begin{lstlisting}[style=CppStyle,caption={Undo state structure}]
struct GameStateData {
    std::vector<std::vector<int>> board;
    int score;
};
\end{lstlisting}

Before each valid move, the engine copies the current board and score. If the move changes the board, the previous state is pushed onto the stack.

\begin{lstlisting}[style=CppStyle,caption={Saving previous state before a valid move}]
GameStateData previousState = {board, score};
bool anyMoved = false;
...
if (anyMoved) {
    history.push(previousState);
    spawnTile();
    if (score > bestScore) bestScore = score;
    checkGameStatus();
}
\end{lstlisting}

\section{Event-Driven Architecture and Game Logic}

The program follows Qt's event-driven structure. It does not use a manual game loop. Instead, Qt calls the appropriate functions when something happens: the user presses a key, a button emits a \texttt{clicked} signal, a message box button is selected, or the hard-mode timer times out.

The user interface calls methods on \texttt{GameEngine}. After each action, \texttt{MainWindow::updateUI} reads the current engine state and redraws labels, colors, button visibility, and dialogs. This one-way synchronization avoids storing duplicate board data in the UI.

\subsection{Keyboard and Button Events}

Keyboard input is handled in \texttt{MainWindow::keyPressEvent}. The arrow keys and WASD keys are mapped to the four slide directions. The \texttt{U} key triggers Undo, while \texttt{R} triggers Restart.

\begin{lstlisting}[style=CppStyle,caption={Keyboard mapping for movement and shortcuts}]
switch (event->key()) {
    case Qt::Key_Up:
    case Qt::Key_W:
        moved = engine.slide(Direction::UP);
        break;
    case Qt::Key_Down:
    case Qt::Key_S:
        moved = engine.slide(Direction::DOWN);
        break;
    case Qt::Key_Left:
    case Qt::Key_A:
        moved = engine.slide(Direction::LEFT);
        break;
    case Qt::Key_Right:
    case Qt::Key_D:
        moved = engine.slide(Direction::RIGHT);
        break;
    case Qt::Key_U:
        onUndoClicked();
        return;
    case Qt::Key_R:
        onRestartClicked();
        return;
}
\end{lstlisting}

If the game has already ended, keyboard movement and Undo are disabled. Only Restart is accepted from the keyboard.

\begin{lstlisting}[style=CppStyle,caption={Input restriction after the game ends}]
if (engine.getState() != GameState::PLAYING) {
    if (event->key() == Qt::Key_R) onRestartClicked();
    return;
}
\end{lstlisting}

The button logic follows the same rules. Restart always starts a new game. Undo does nothing if the game is not in the \texttt{PLAYING} state. Mode buttons switch modes without resetting the board or score.

\begin{lstlisting}[style=CppStyle,caption={Normal Mode switch preserves the board}]
void MainWindow::setModeNormal()
{
    if (engine.getMode() == GameMode::NORMAL ||
        engine.getState() != GameState::PLAYING) {
        return;
    }
    engine.setMode(GameMode::NORMAL);
    hardModeTimer->stop();
    updateUI();
}
\end{lstlisting}

Hard Mode starts the timer when selected. Leaving Hard Mode stops the timer. This prevents random moves from continuing after the player returns to Normal or Unlimited Mode.

\begin{lstlisting}[style=CppStyle,caption={Hard Mode starts the timer without restarting}]
void MainWindow::setModeHard()
{
    if (engine.getMode() == GameMode::HARD ||
        engine.getState() != GameState::PLAYING) {
        return;
    }
    engine.setMode(GameMode::HARD);
    resetHardModeTimer();
    updateUI();
}
\end{lstlisting}

\subsection{Signals, Slots, and Connections}

Qt's signal-slot mechanism connects UI events to functions in \texttt{MainWindow}. The connections are created in \texttt{setupUI} for buttons and in the constructor for the timer. This keeps the UI reactive: buttons do not need to be checked manually, and the timer event is delivered by Qt.

\begin{longtable}{|p{4.4cm}|p{4.2cm}|p{5.4cm}|}
\hline
\textbf{Signal} & \textbf{Slot / Function} & \textbf{Purpose} \\
\hline
\texttt{normalModeBtn.clicked()} & \texttt{setModeNormal()} & Switch to Normal Mode while preserving board and score. \\
\hline
\texttt{unlimitedModeBtn.clicked()} & \texttt{setModeUnlimited()} & Switch to Unlimited Mode or continue after a win. \\
\hline
\texttt{hardModeBtn.clicked()} & \texttt{setModeHard()} & Switch to Hard Mode and start the five-second timer. \\
\hline
\texttt{undoBtn.clicked()} & \texttt{onUndoClicked()} & Restore the previous board and score during active play. \\
\hline
\texttt{restartBtn.clicked()} & \texttt{onRestartClicked()} & Start a new game, reset score, and clear undo history. \\
\hline
\texttt{hardModeTimer.timeout()} & \texttt{onHardModeTimeout()} & Perform a random valid move after five seconds in Hard Mode. \\
\hline
\caption{Main Qt signal-slot connections used in the project.}
\end{longtable}

\subsection{Move and Merge Logic}

The core movement logic is in \texttt{GameEngine::slide} and \texttt{GameEngine::slideArray}. The \texttt{slide} function handles direction-specific row and column extraction. The \texttt{slideArray} function handles the actual compression and merging of a one-dimensional line.

For left and right moves, each board row is copied into a temporary vector. For right moves, the vector is reversed before merging, then reversed back before being written to the board. For up and down moves, the same idea is applied to columns. This keeps one merge algorithm consistent across all four directions.

\begin{lstlisting}[style=CppStyle,caption={Using one-dimensional lines for row movement}]
if (dir == Direction::LEFT || dir == Direction::RIGHT) {
    for (int r = 0; r < BOARD_N; ++r) {
        std::vector<int> line = board[r];
        if (dir == Direction::RIGHT) std::reverse(line.begin(), line.end());

        if (slideArray(line)) anyMoved = true;

        if (dir == Direction::RIGHT) std::reverse(line.begin(), line.end());
        board[r] = line;
    }
}
\end{lstlisting}

The merge algorithm has three phases:

\begin{enumerate}
    \item Copy non-zero values into a new line.
    \item Scan adjacent values from the movement side and merge equal pairs.
    \item Append zeros until the line returns to its original size.
\end{enumerate}

\begin{lstlisting}[style=CppStyle,caption={Core merge algorithm}]
for (int val : line) {
    if (val != 0) newLine.push_back(val);
}

for (size_t i = 0; i + 1 < newLine.size(); ++i) {
    if (newLine[i] == newLine[i+1]) {
        newLine[i] *= 2;
        score += newLine[i];
        newLine.erase(newLine.begin() + i + 1);
        moved = true;
    }
}

while (newLine.size() < line.size()) {
    newLine.push_back(0);
}
\end{lstlisting}

This approach also handles the required merge order. If three equal tiles appear in a row, the pair closer to the movement border merges first. For example, moving left transforms \texttt{[2, 2, 2, 0]} into \texttt{[4, 2, 0, 0]}. Moving right is handled by reversing the line first, so the rightmost pair merges first.

Invalid moves are detected by comparing the processed line with the original line and by tracking whether any merge occurred. If the whole board is unchanged after a requested move, no new tile is generated and the previous state is not pushed to the undo stack.

\subsection{Timer-Based Events}

Hard Mode uses a \texttt{QTimer}. When Hard Mode begins, \texttt{resetHardModeTimer} starts the timer with a five-second interval. The timer is also restarted after each valid player move and after each valid automatic move.

\begin{lstlisting}[style=CppStyle,caption={Starting the Hard Mode timer}]
void MainWindow::resetHardModeTimer()
{
    hardModeTimer->start(5000);
}
\end{lstlisting}

When the timer expires, the program attempts random directions until a valid move is made. The automatic move uses the same \texttt{engine.slide} function as player input. This is important because it means timer-based moves update score, undo history, tile generation, and win/loss state consistently.

\begin{lstlisting}[style=CppStyle,caption={Hard Mode timeout logic}]
void MainWindow::onHardModeTimeout()
{
    std::vector<Direction> dirs =
        {Direction::UP, Direction::DOWN, Direction::LEFT, Direction::RIGHT};
    bool moved = false;

    while (!moved && engine.getState() == GameState::PLAYING) {
        int r = std::rand() % 4;
        moved = engine.slide(dirs[r]);
    }

    if (moved) {
        updateUI();
        resetHardModeTimer();
    }
}
\end{lstlisting}

\section{Code Structure}

The project uses a small number of source files, each with a clear role.

\begin{table}[H]
\centering
\begin{tabular}{|p{4cm}|p{10cm}|}
\hline
\textbf{File} & \textbf{Responsibility} \\
\hline
\texttt{main.cpp} & Creates the Qt application object, creates the main window, shows it, and starts the event loop. \\
\hline
\texttt{mainwindow.h/cpp} & Defines and implements the user interface, keyboard handling, button slots, timer behavior, message boxes, tile colors, and UI synchronization. \\
\hline
\texttt{gameengine.h/cpp} & Defines and implements the board state, movement rules, merge rules, tile generation, score tracking, best score tracking, undo stack, win detection, and loss detection. \\
\hline
\texttt{CMakeLists.txt} & Defines the CMake project, Qt dependency, source files, executable target, and linking configuration. \\
\hline
\end{tabular}
\caption{Main files and responsibilities.}
\end{table}

The most important design decision is the separation between UI and game logic. \texttt{GameEngine} does not know about labels, buttons, layouts, or Qt events. It only represents the rules of the game. \texttt{MainWindow} owns the visual widgets and calls the engine through a small public interface such as \texttt{restart}, \texttt{slide}, \texttt{undo}, \texttt{getTile}, and \texttt{setMode}. This interface is enough for \texttt{MainWindow} to display the board, execute player actions, and update controls without directly modifying private engine fields.

\subsection*{UI Synchronization}

\texttt{MainWindow::updateUI} is called after every operation that can change the board or state. It reads each tile value from the engine and updates the corresponding label. It also updates scores, button states, tile sizes, and end-game dialogs.

\begin{lstlisting}[style=CppStyle,caption={Synchronizing tile labels with the engine board}]
for (int r = 0; r < BOARD_N; ++r) {
    for (int c = 0; c < BOARD_M; ++c) {
        int val = engine.getTile(r, c);
        if (val == 0) {
            boardLabels[r][c]->setText("");
            boardLabels[r][c]->setStyleSheet(
                "background-color: #cdc1b4; border-radius: 5px;");
        } else {
            boardLabels[r][c]->setText(QString::number(val));
            boardLabels[r][c]->setStyleSheet(getTileColor(val));
        }
    }
}
\end{lstlisting}

\subsection*{End-Game Control State}

The UI control state is updated from the engine state. During active play, mode and undo controls are visible. After a loss, only Restart remains visible. After a win, the player can restart or continue in Unlimited Mode.

\begin{lstlisting}[style=CppStyle,caption={Control visibility and availability}]
const bool playing = engine.getState() == GameState::PLAYING;
const bool won = engine.getState() == GameState::WON;

normalModeBtn->setVisible(playing);
unlimitedModeBtn->setVisible(playing || won);
hardModeBtn->setVisible(playing);
undoBtn->setVisible(playing);

normalModeBtn->setEnabled(playing && engine.getMode() != GameMode::NORMAL);
unlimitedModeBtn->setEnabled((playing && engine.getMode() != GameMode::UNLIMITED)
                             || won);
hardModeBtn->setEnabled(playing && engine.getMode() != GameMode::HARD);
undoBtn->setEnabled(playing);
restartBtn->setEnabled(true);
\end{lstlisting}

\section{Testing and Debugging}

The implementation was checked against the important cases described in the specification and in project clarifications.

\begin{itemize}
    \item \textbf{Initial state:} Restart creates an empty board and spawns exactly two tiles with value 2 in random empty positions.
    \item \textbf{Invalid move:} If a move does not change the board, the move returns false, no undo state is saved, and no new tile is generated.
    \item \textbf{Repeated equal tiles:} Lines with multiple equal tiles merge from the movement side. This satisfies the rule that the pair closer to the border merges first.
    \item \textbf{Full board with moves available:} The game does not immediately end when the board is full. The engine checks adjacent equal tiles before deciding that the player lost.
    \item \textbf{Full board without moves:} If there are no empty cells and no adjacent equal tiles, the engine changes the state to \texttt{LOST}.
    \item \textbf{Undo to the beginning:} Undo is safe even if the history stack is empty. Pressing Undo more times than the number of valid moves has no effect and does not crash the game.
    \item \textbf{Best score with Undo:} Undo restores the current score but not the best score.
    \item \textbf{Mode switching:} Switching modes preserves board and score instead of restarting the game.
    \item \textbf{Win condition:} Reaching the target tile in Normal or Hard Mode shows the win message. The player can restart or continue in Unlimited Mode.
    \item \textbf{Unlimited Mode:} Reaching the target tile does not end the game, and the player can continue to larger values.
    \item \textbf{Hard Mode:} If the timer expires, the game executes a random valid move and then restarts the timer.
    \item \textbf{Responsive board:} Tile size is recalculated so that rectangular boards such as 8 by 6 or 10 by 6 still use square tiles.
\end{itemize}

The loss condition is implemented by \texttt{canMove}. It first checks for empty cells. If none exist, it checks horizontally and vertically for adjacent equal tiles. Only if both checks fail does it return false.

\begin{lstlisting}[style=CppStyle,caption={Loss-condition helper}]
bool GameEngine::canMove() const {
    for (int r = 0; r < BOARD_N; ++r) {
        for (int c = 0; c < BOARD_M; ++c) {
            if (board[r][c] == 0) return true;
        }
    }

    for (int r = 0; r < BOARD_N; ++r) {
        for (int c = 0; c < BOARD_M - 1; ++c) {
            if (board[r][c] == board[r][c+1]) return true;
        }
    }

    for (int r = 0; r < BOARD_N - 1; ++r) {
        for (int c = 0; c < BOARD_M; ++c) {
            if (board[r][c] == board[r+1][c]) return true;
        }
    }

    return false;
}
\end{lstlisting}

This structure avoids a common bug where a full board is treated as lost even though a merge is still possible.

\section{Challenges and Solutions}

\textbf{Preventing illegal double merges.} The most important mechanical challenge in 2048 is preventing one tile from merging twice in a single move. The solution is to process a compacted one-dimensional line and remove the second tile of a pair immediately after merging. Because the consumed tile is erased, it cannot be merged again in that move.

\textbf{Correct merge ordering.} When three or more equal tiles appear in the same line, the pair closer to the movement border must merge first. The implementation solves this by always scanning the temporary line from the movement side. For right and down moves, the line is reversed before processing, so the same left-to-right scan still represents movement toward the border.

\textbf{Keeping movement logic maintainable.} Instead of writing four separate algorithms for left, right, up, and down, the code converts rows and columns into vectors and reuses \texttt{slideArray}. This reduces duplication and helps keep scoring, invalid move detection, and merge behavior consistent across directions.

\textbf{Synchronizing UI and state.} A Qt interface can easily become inconsistent if UI labels store their own separate version of the board. The implementation avoids this by making \texttt{GameEngine} the single source of truth. After every event, \texttt{updateUI} reads from the engine and redraws the visible board.

\textbf{Handling mode switching.} Mode buttons should change the rules without starting a new game. This required separating mode changes from restart behavior. Hard Mode also needed special handling: entering the mode starts the timer, while leaving it stops the timer.

\textbf{Disabling input after game end.} The specification requires movement, Undo, and Hard Mode random moves to be disabled after the game ends. The solution is split between engine and UI. The engine rejects slides and undo actions unless the state is \texttt{PLAYING}, while the UI hides or disables controls and ignores movement keys after win or loss.

\textbf{Supporting configurable board sizes.} A fixed tile size works for the default 4 by 4 board but breaks down for larger rectangular boards. The solution is to compute a square tile size from the available width, available height, number of rows, and number of columns. This keeps the UI usable when constants such as \texttt{N} and \texttt{M} are changed.

\section{Conclusion}

The project implements a complete Qt-based version of 2048 with Normal, Unlimited, and Hard modes. The game includes movement, merging, scoring, best score tracking, unlimited Undo, Restart, win handling, loss handling, configurable constants, and a responsive tile grid. The executable is configured through CMake with the target name \texttt{2048}, and the repository keeps source and configuration files separate from generated build output.

The strongest part of the design is the separation between \texttt{GameEngine} and \texttt{MainWindow}. The engine owns the game rules and state, while the UI owns widgets and event handling. This separation makes it easier to understand the project during a demo session because movement rules can be explained independently from Qt layout code.

The most difficult implementation problems were the movement/merge algorithm, preserving correct behavior across all four directions, and handling mode changes without accidentally restarting the game. The current implementation solves these by using one-dimensional line processing, direction reversal, and explicit mode-switching functions.

Possible future improvements include adding smooth tile animations, showing a visible countdown in Hard Mode, and replacing message boxes with a custom overlay widget. These improvements would polish the presentation, but the current version already provides the required gameplay behavior and event-driven structure.

\section*{AI Usage Statement}

AI assistance was used during the preparation of this project for limited support tasks such as reviewing the report structure, improving wording, and checking whether the written explanations matched the source code. The implementation, design decisions, and final technical content were implemented by the authors but only reviewed by AI models. The authors understand the code structure, game logic, Qt event handling, and the behavior described in this report.

\end{document}
